\begin{theorem}[Cohen structure theorem]
\label{theorem-cohen-structure-theorem}
Let $(R, \mathfrak m)$ be a complete local ring.
\begin{enumerate}
\item $R$ has a coefficient ring (see
Definition \ref{definition-coefficient-ring}),
\item if $\mathfrak m$ is a finitely generated ideal, then
$R$ is isomorphic to a quotient
$$
\Lambda[[x_1, \ldots, x_n]]/I
$$
where $\Lambda$ is either a field or a Cohen ring.
\end{enumerate}
\end{theorem}

\begin{proof}
Let us prove a coefficient ring exists.
First we prove this in case the characteristic of the residue field $\kappa$
is zero. Namely, in this case we will prove by induction
on $n > 0$ that there exists a section
$$
\varphi_n : \kappa \longrightarrow R/\mathfrak m^n
$$
to the canonical map $R/\mathfrak m^n \to \kappa = R/\mathfrak m$.
This is trivial for $n = 1$. If $n > 1$, let $\varphi_{n - 1}$ be given.
The field extension $\kappa/\mathbf{Q}$ is formally smooth by
Proposition \ref{proposition-characterize-separable-field-extensions}.
Hence we can find the dotted arrow
in the following diagram
$$
\xymatrix{
R/\mathfrak m^{n - 1} &
R/\mathfrak m^n \ar[l] \\
\kappa \ar[u]^{\varphi_{n - 1}} \ar@{..>}[ru] & \mathbf{Q} \ar[l] \ar[u]
}
$$
This proves the induction step. Putting these maps together
$$
\lim_n\ \varphi_n : \kappa \longrightarrow
R = \lim_n\ R/\mathfrak m^n
$$
gives a map whose image is the desired coefficient ring.

\medskip\noindent
Next, we prove the existence of a coefficient ring in the case
where the characteristic of the residue field $\kappa$ is $p > 0$.
Namely, choose a Cohen ring $\Lambda$ with $\kappa = \Lambda/p\Lambda$,
see Lemma \ref{lemma-cohen-rings-exist}. In this case we will prove by
induction on $n > 0$ that there exists a map
$$
\varphi_n :
\Lambda/p^n\Lambda
\longrightarrow
R/\mathfrak m^n
$$
whose composition with the reduction map $R/\mathfrak m^n \to \kappa$
produces the given isomorphism $\Lambda/p\Lambda = \kappa$. This is trivial
for $n = 1$. If $n > 1$, let $\varphi_{n - 1}$ be given.
The ring map $\mathbf{Z}/p^n\mathbf{Z} \to \Lambda/p^n\Lambda$
is formally smooth by Lemma \ref{lemma-cohen-ring-formally-smooth}.
Hence we can find the dotted arrow
in the following diagram
$$
\xymatrix{
R/\mathfrak m^{n - 1} &
R/\mathfrak m^n \ar[l] \\
\Lambda/p^n\Lambda \ar[u]^{\varphi_{n - 1}} \ar@{..>}[ru] &
\mathbf{Z}/p^n\mathbf{Z} \ar[l] \ar[u]
}
$$
This proves the induction step. Putting these maps together
$$
\lim_n\ \varphi_n :
\Lambda = \lim_n\ \Lambda/p^n\Lambda
\longrightarrow
R = \lim_n\ R/\mathfrak m^n
$$
gives a map whose image is the desired coefficient ring.

\medskip\noindent
The final statement of the theorem follows readily. Namely, if
$y_1, \ldots, y_n$ are generators of the ideal $\mathfrak m$,
then we can use the map $\Lambda \to R$ just constructed
to get a map
$$
\Lambda[[x_1, \ldots, x_n]] \longrightarrow R,
\quad x_i \longmapsto y_i.
$$
Since both sides are $(x_1, \ldots, x_n)$-adically complete
this map is surjective by Lemma \ref{lemma-completion-generalities}
as it is surjective modulo $(x_1, \ldots, x_n)$ by
construction.
\end{proof}